\scanpage{9}
\uline{Axiomes} $1^\circ:6^\circ$, et.

A) Si $S'$ local noeth complet, \inserted{[de dim $1$ si on veut]} la promodularité de
\[
 \unsure{X|Z}\xrightarrow{\unsure{\eta|y}}F_{S'}
\]
en $z$ sur le pt fermé de $S'$ ($\unsure{X|Z}$ fini sur $S'$) implique $\unsure{\eta|y}$ \deleted{modulaire} imm. ouverte (et on peut supposer $F|(S'-\{s'\})$ repr.)

B) \deleted{Si \unsure{aux des}} $S'$ noeth sur $S$, $\unsure{X|Z}$ sur $S$, $\unsure{\eta|y}:\unsure{X|Z}\mathrel{\unsure{-|\longrightarrow}}F_{S'}$, et $x\in X$ sont tels que $\unsure{\eta|y}$ \unsure{modulaire} en $x$, alors $\exists\ U\ni x$ voisinage ouvert, t.q. $\unsure{\eta|y}|U:U\mathrel{\unsure{-|\longrightarrow}}F_{S'}$ soit une imm. ouverte [On \unsure{peut|pt} $\unsure{\simeq|\text{même}}$ supposer que $X\times_{S'}\mathcal O_{S',s'}\longrightarrow F_{\operatorname{Spec}\mathcal O_{S',s'}}$ est une imm. ouverte]

\uuline{Cas particuliers}

\deleted{(i) \illegible{1} $F\longrightarrow S$ un \uline{monomorphisme}. Alors (A) et (B) sont contenus dans}

\deleted{(C) Pour $\forall\ S'$ noethérien$/S$, \unsure{lisse}}
\inserted{$X\longrightarrow S'$ un \deleted{\unsure{mor.}} \unsure{mor.}}
\deleted{et $\unsure{\eta|y}:X\mathrel{\unsure{-|\longrightarrow}}F_{S'}$, l'ens des $x\in X$ en lesquels $\unsure{\eta|y}$ est \inserted{inf.} étale est ouvert.}

\margin{\uuline{Lemme}\quad On peut supposer en A resp. B que $X\longrightarrow S'$ une \uuline{imm.} \uline{fermée}.}
\margin{définie \unsure{si} on veut par un idéal de carré nul}

Si $F\mathrel{\unsure{-|\longrightarrow}}S$ \uline{non ramifié}, les conditions (A) et (B), sont \inserted{alors} contenues dans

(C) Si $S'$ noeth. sur $S$, \deleted{$S'_0$} \inserted{$S'$ \uline{imm. fermée}, définie par un idéal de carré nul} \deleted{\unsure{tf.}} sur $S'$, $\unsure{\eta|y}:X\mathrel{\unsure{-|\longrightarrow}}F_{S'}$, l'ens \inserted{$U$} des pts en lesquels $\unsure{\eta|y}$ est \uline{inf\textsuperscript{t}.} \uline{étale} est \uline{ouvert} i.e.

$(C_1)$ stable par générisation [est \deleted{\unsure{un}} équivalent à A]

$(C_2)$ si $x\in U$, alors \deleted{l'ens} $U\cap\bar x$ est un voisinage de $x$ dans $\bar x$.
\margin{\figure{\begin{tikzpicture}\draw (0,2)--(0,3);\draw (0,0)--(0,.5);\end{tikzpicture}}}

\uline{N.B.} La condition $(C_1)$ est automatique si $F\longrightarrow S$ est \uline{étale}
\mspage{8}
